\documentclass[10pt,a4paper]{amsart}
\usepackage[margin=19mm]{geometry}
\usepackage{amsmath,amssymb,mathtools}
\usepackage[scr=rsfs,cal=euler]{mathalfa}
\usepackage{xfrac}
\usepackage[unicode,hidelinks,bookmarks=false]{hyperref}
\newcommand{\ie}{i.e.\ }
\newcommand{\cf}{cf.\ }
\newcommand{\resp}{respectively\ }
\newcommand{\Subsection}{\subsection}
\providecommand{\nobreakdash}{\nobreak\hbox{-}}
\setlength{\emergencystretch}{2em}
\multlinegap0pt
\usepackage{amssymb}
\usepackage{mathtools}
\usepackage[shortlabels]{enumitem}
\usepackage{tcolorbox}
\usepackage{float}
\usepackage{xcolor}
\usepackage{tensor}
\newtheorem{lemma}{Lemma}
\title{Dirichlet Beta Analogues, Sign Alternation, and Lerch Unification of Hyperbolic and Logarithmic Tangent Integrals}
\author{Luc Rams\`es Talla Waffo}
\date{Source: arXiv:2608.27717v1 (CC BY 4.0)}
\begin{document}
\maketitle

\noindent\emph{Source and licence.} Dirichlet Beta Analogues, Sign Alternation, and Lerch Unification of Hyperbolic and Logarithmic Tangent Integrals. Version arXiv:2608.27717v1 (CC BY 4.0), distributed by arXiv under the Creative Commons Attribution 4.0 International licence, http://creativecommons.org/licenses/by/4.0/, as recorded in the arXiv licence metadata for this exact version and shown on the abstract page https://arxiv.org/abs/2608.27717v1. Full licence notice: \texttt{licenses/arxiv-2608.27717-CC-BY-4.0.txt}.\par\smallskip

\noindent\textit{Editorial scope.} Counted unit: the article's lemma lemma (source0.tex lines 372--400). The article's complete original proof of it is delivered with it (source0.tex lines 402--420, one \texttt{proof} environment of 19 lines). No mathematics is written, repaired or re-derived: the statement and the proof are the article's own lines, in the article's own notation, and no retained line is substituted or re-flowed.  No further source-local context is needed: every cross-reference of every retained line resolves inside the two delivered spans.  Nothing was omitted from any retained span, so the package carries no omission record.  No retained line contains a citation, so the wrapper carries no bibliography and no bibliography entry.  No retained line contains a figure environment, an image-inclusion command or drawing code, so no image is delivered and the package declares no asset dependency.  Editorial formatting only, in addition: an AMS \texttt{amsart} preamble that restates the article's own declarations and macros used by the retained lines, namely the environments \texttt{lemma} on separate counters, together with the standard packages those lines need.  The retained environments print in this document as Lemma 1 in order of appearance, whereas the article prints the same items with the numbering of its own full document.  The complete native source of the article as distributed in the arXiv e-print for this exact version is bundled unchanged as \texttt{sources/208734/source0.tex}.
\begin{lemma}[Derivative-polynomial structure]
\label[lemma]{lem:sech-tanh-derivative-polynomial}
Let \(m\geq1\) be an integer and define
\(f_m(x):=\tanh^m x/\cosh x\).
For every integer \(r\geq0\), there exists a polynomial \(P_{m,r}\)
such that
\[
f_m^{(r)}(x)
=
\frac{1}{\cosh x}P_{m,r}(\tanh x).
\]
The polynomials satisfy \(P_{m,0}(t)=t^m\) and
\[
P_{m,r+1}(t)
=
(1-t^2)P_{m,r}'(t)-tP_{m,r}(t).
\]
If \(0\leq r\leq m\), then the smallest and largest powers occurring
in \(P_{m,r}\) are \(t^{m-r}\) and \(t^{m+r}\), respectively,
and only powers of the same parity occur. In particular, if \(m+r\)
is odd, then
\[
P_{m,r}(t)
=
\sum_{k=(m-r-1)/2}^{(m+r-1)/2}
a_{m,r,k}\,t^{2k+1}
\]
for suitable coefficients \(a_{m,r,k}\).
\end{lemma}

\begin{proof}
Put \(t=\tanh x\). Since \(dt/dx=1-t^2\) and
\(\dfrac{d}{dx}(\cosh^{-1}x)=-t\cosh^{-1}x\), we obtain
\[
\frac{d}{dx}
\left(\frac{P(t)}{\cosh x}\right)
=
\frac{(1-t^2)P'(t)-tP(t)}{\cosh x},
\]
which proves the recurrence.

Starting from \(P_{m,0}(t)=t^m\), induction shows that
\(P_{m,r}\) contains only powers
\(t^{m-r},t^{m-r+2},\ldots,t^{m+r}\).
At each step the coefficient of the lowest power is multiplied by
\(m-r\), while that of the highest power is multiplied by
\(-(m+r+1)\). Thus both endpoint coefficients are nonzero for
\(0\leq r\leq m\).
\end{proof}

\end{document}
